%% kapa: exam=1 round=28 year=2017 period=2 form=A subject=appraisal_law subject_ko=감정평가관계법규 seq=35 number=35
%% source: Q-Net 감정평가사 기출문제 (cst00309) artlSeq=5199603 file=extracted/1008/5199603/2017년 제28회 감정평가사 1차시험 문제지/2교시 A형.pdf
%% conversion: kordoc 4.24.51 -> tools/gen1.py (auto); review=proofread-sonnet; flags=none
\begin{question}{35}
부동산등기법령상 ‘이의’에 관한 설명으로 옳은 것을 모두 고른 것은?
\begin{framedtext}
ㄱ. 이의의 신청은 등기소에 이의신청서를 제출하는 방법으로 한다. \par
ㄴ. 새로운 사실이나 새로운 증거방법을 근거로 이의신청을 할 수 있다. \par
ㄷ. 이의에는 집행정지의 효력이 없다. \par
ㄹ. 등기관은 이의가 이유 없다고 인정하면 이의신청일부터 7일 이내에 의견을 붙여 이의신청서를 관할 지방법원에 보내야 한다. \par
\end{framedtext}
\begin{choices}
  \item ㄱ, ㄴ
  \item ㄱ, ㄷ
  \item ㄴ, ㄷ
  \item ㄴ, ㄹ
  \item ㄷ, ㄹ
\end{choices}
\end{question}
\answer{2}
